% 第三章 混合权重评估模型构建
% ch3-model.tex

\subsection{问题数学化建模}

\subsubsection{多准则决策问题定义}

奥运项目评估可抽象为多准则决策问题。设：
\begin{itemize}
\item 备选项目集 $S = \{s_1, s_2, \ldots, s_n\}$
\item 评价指标集 $C = \{c_1, c_2, \ldots, c_6\}$（六维指标）
\item 决策矩阵 $X = [x_{ij}]_{n \times 6}$
\end{itemize}

目标：确定权重向量 $W = (w_1, w_2, \ldots, w_6)$，使综合评分合理反映项目优劣。

\subsubsection{综合评分模型}

\begin{equation}
v_i = \sum_{j=1}^{6} w_j \cdot E_{ij}
\label{eq:score}
\end{equation}
其中 $E_{ij}$ 为项目 $i$ 在维度 $j$ 上的综合值（由子指标计算）。

\subsection{六维评价指标体系构建}

基于IOC评估标准和Song等\cite{song2025quantifying}的研究，构建六维评价指标体系。

\begin{table}[htbp]
\centering
\caption{六维评价指标体系}
\label{tab:indicators}
\begin{tabular}{cll}
\toprule
编号 & 指标名称 & 说明 \\
\midrule
$c_1$ & 流行度 & 观众收视率、参与人数、社交媒体讨论度 \\
$c_2$ & 性别平等 & 男女参赛比例、混合项目数量 \\
$c_3$ & 可持续性 & 碳排放、场馆成本、资源消耗 \\
$c_4$ & 包容性 & 残奥关联、发展中国家参与度 \\
$c_5$ & 创新性 & 年轻化程度、技术革新、新近加入 \\
$c_6$ & 安全性 & 受伤率、医疗保障需求 \\
\bottomrule
\end{tabular}
\end{table}

\subsection{AHP主观权重计算模型}

\subsubsection{判断矩阵构建}

邀请$k$位专家对六维指标进行两两比较，得到判断矩阵：
\begin{equation}
A = \begin{bmatrix}
1 & a_{12} & a_{13} & a_{14} & a_{15} & a_{16} \\
1/a_{12} & 1 & a_{23} & a_{24} & a_{25} & a_{26} \\
1/a_{13} & 1/a_{23} & 1 & a_{34} & a_{35} & a_{36} \\
1/a_{14} & 1/a_{24} & 1/a_{34} & 1 & a_{45} & a_{46} \\
1/a_{15} & 1/a_{25} & 1/a_{35} & 1/a_{45} & 1 & a_{56} \\
1/a_{16} & 1/a_{26} & 1/a_{36} & 1/a_{46} & 1/a_{56} & 1
\end{bmatrix}
\label{eq:jm}
\end{equation}

\subsubsection{多专家矩阵整合}

采用几何平均法整合$k$位专家的判断矩阵：
\begin{equation}
\bar{a}_{ij} = \left( \prod_{l=1}^{k} a_{ij}^{(l)} \right)^{1/k}
\label{eq:geomean}
\end{equation}

\subsubsection{权重计算与一致性检验}

求解特征值问题$AW = \lambda_{\max} W$，归一化得权重向量$W^{AHP}$。

一致性检验：$CR < 0.1$时判断矩阵有效。

\subsubsection{模拟专家权重结果}

基于模拟专家访谈（见附录），得到AHP主观权重：
\begin{table}[htbp]
\centering
\caption{AHP主观权重（模拟数据）}
\label{tab:ahp_weights}
\begin{tabular}{ccc}
\toprule
指标 & 权重 & 排名 \\
\midrule
流行度 & 0.360 & 1 \\
性别平等 & 0.176 & 2 \\
安全性 & 0.174 & 3 \\
可持续性 & 0.112 & 4 \\
包容性 & 0.095 & 5 \\
创新性 & 0.063 & 6 \\
\bottomrule
\end{tabular}
\end{table}

一致性比率 $CR = 0.0067 < 0.1$，通过检验。

\subsection{EWM客观权重计算模型}

\subsubsection{数据标准化}

正向指标（越大越好）：
\begin{equation}
\tilde{x}_{ij} = \frac{x_{ij} - \min_i x_{ij}}{\max_i x_{ij} - \min_i x_{ij}}
\label{eq:pos_norm}
\end{equation}

负向指标（越小越好）：
\begin{equation}
\tilde{x}_{ij} = \frac{\max_i x_{ij} - x_{ij}}{\max_i x_{ij} - \min_i x_{ij}}
\label{eq:neg_norm}
\end{equation}

\subsubsection{信息熵与权重计算}

见第\ref{sec:literature}章公式(\ref{eq:info_entropy})-(\ref{eq:entropy_weight})。

\subsubsection{两层EWM权重计算}

本研究采用两层结构：

\textbf{第一层}：六维指标权重（与AHP混合）
\textbf{第二层}：各维度子指标权重（纯EWM）

\subsection{混合权重组合策略}

\subsubsection{线性组合模型}

\begin{equation}
W = \alpha \cdot W^{AHP} + (1 - \alpha) \cdot W^{EWM}
\label{eq:hybrid_weight}
\end{equation}

\subsubsection{混合权重的数学性质}

\textbf{性质1（有界性）}：
\begin{equation}
\min(w_j^{AHP}, w_j^{EWM}) \leq w_j \leq \max(w_j^{AHP}, w_j^{EWM})
\end{equation}

\textbf{性质2（单调性）}：
\begin{equation}
\frac{\partial w_j}{\partial \alpha} = w_j^{AHP} - w_j^{EWM}
\end{equation}

\textbf{性质3（凸组合）}：混合权重是$W^{AHP}$和$W^{EWM}$的凸组合，满足归一化条件。

\subsubsection{组合系数确定}

推荐$\alpha = 0.5$（主客观权重等权），并通过灵敏度分析验证。

\subsection{模型有效性分析}

\subsubsection{与传统方法对比}

\begin{table}[htbp]
\centering
\caption{权重确定方法对比}
\label{tab:weight_comparison}
\begin{tabular}{lccc}
\toprule
方法 & 主观性 & 客观性 & 优势 \\
\midrule
纯AHP & 高 & 低 & 融入专家知识 \\
纯EWM & 低 & 高 & 数据驱动 \\
混合模型 & 中 & 中 & 综合主客观 \\
\bottomrule
\end{tabular}
\end{table}

\subsubsection{模型优势}

\begin{enumerate}
\item 兼顾主观判断与客观数据
\item 两层结构提高灵活性
\item 完整的数学推导支撑
\end{enumerate}

\subsection{本章小结}

本章构建了基于AHP-EWM的混合权重评估模型。首先将奥运项目评估问题抽象为多准则决策问题，建立六维评价指标体系。然后分别推导了AHP主观权重和EWM客观权重的计算方法。最后提出线性组合的混合权重模型，分析了其数学性质。下一章将通过实验验证模型的有效性。